\section{Introduction}
%================================================
\begin{frame}{Active Matter Systems}

	\begin{columns}[T]

		%-----------------------------------
		\column{0.38\textwidth}

		\begin{itemize}
			\item Active matter consists of entities that consume energy and convert it into motion.

			\item Active systems are found in both living and synthetic systems.

			\item These systems exhibit persistence, collective behaviour, and non-equilibrium transport.
		\end{itemize}

		\vspace{-0.2cm}

		\begin{alertblock}{Motivating Question}
			What mechanism allows active droplets to remember where they have been?
		\end{alertblock}

		%-----------------------------------
		\column{0.58\textwidth}

		\begin{block}{Examples of Active Matter}

			\centering

			\begin{tabular}{ccc}

				\includegraphics[height=1.9cm]{Starling_murmuration.jpg}
				 &
				\includegraphics[height=1.9cm]{bacteria.png}
				 &
				\includegraphics[height=1.9cm]{lab_droplet_image.png}

				\\[0.05cm]

				{\small Bird flock}
				 &
				{\small Bacteria}
				 &
				{\small Active droplet}
			\end{tabular}

		\end{block}
		\vspace{0.1cm}
		{\tiny
			\textbf{Image Sources}
			\begin{itemize}
				\item Bird flock: \emph{A murmuration of starlings at Gretna}, adapted from Walter Baxter (2011), Wikimedia Commons.

				\item Bacterial colony: Growing \textit{E. coli} microcolony exhibiting active-nematic ordering, adapted from Dell'Arciprete \textit{et al.}, \emph{Nature Communications} (2018).

				\item Active droplet: Self-propelled Belousov--Zhabotinsky (BZ) droplet driven by chemical oscillations and Marangoni flow, Nonlinear Dynamics Laboratory, NITK.
			\end{itemize}
		}
	\end{columns}


\end{frame}
%================================================

%================================================
\begin{frame}{Why Active Droplets?}

	\begin{columns}[T]

		%-----------------------------------
		\column{0.48\textwidth}

		\centering

		\includegraphics[width=\linewidth]{BZ_active_droplet.png}

		\vspace{0.15cm}

		{\scriptsize
			BZ active droplet exhibiting internal chemical-wave dynamics.
			Adapted from Meshram and Shajahan (2025).
		}

		%-----------------------------------
		\column{0.48\textwidth}

		\begin{block}{Key Features}

			\begin{itemize}

				\item Simple and controllable active matter system.

				\item Self-propelled through internally generated chemical activity.

				\item Continuously modifies its surrounding environment.

				\item Provides a natural platform for studying memory and feedback effects.

			\end{itemize}

		\end{block}

	\end{columns}

	\vspace{0.2cm}

	\begin{alertblock}{Hypothesis}
		The self-generated chemical field may retain information about previously visited locations and influence future motion.
	\end{alertblock}

\end{frame}
%================================================

%------------------------------------------------
\begin{frame}{Observation: Self-Avoiding Motion}

	\begin{columns}[T]

		\column{0.60\textwidth}

		\centering
		\includegraphics[height=0.62\textheight]{for_2.1_droplets_droplet_3_trajectory.png}

		\column{0.40\textwidth}

		\begin{block}{Observation}
			The droplet tends to avoid regions it has previously visited.
		\end{block}

		\vspace{0.2cm}

		\begin{itemize}
			\item Motion is not purely random.
			\item Trajectories exhibit self-avoidance.
			\item Past motion appears to influence future motion.
		\end{itemize}

	\end{columns}

	\vspace{0.15cm}

	\begin{alertblock}{Research Question}
		What mechanism allows the droplet to remember where it has been?
	\end{alertblock}

\end{frame}
%------------------------------------------------
\begin{frame}{A Possible Explanation: Memory}

	\centering

	\begin{tikzpicture}[
			node distance=2.5cm,
			every node/.style={
					draw,
					rounded corners,
					align=center,
					minimum width=3.8cm,
					minimum height=1cm
				},
			>=Stealth
		]

		\node (motion) {Droplet Motion};

		\node (deposit) [right=of motion]
		{Chemical\\Deposition};

		\node (memory) [below=of deposit,
			fill=red!15,
			draw=red!70!black]
		{Chemical Field\\(Memory Storage)};

		\node (retrieve) [left=of memory,
			fill=red!10]
		{Gradient Sensing\\(Memory Retrieval)};

		\draw[->,thick] (motion) -- (deposit);
		\draw[->,thick] (deposit) -- (memory);
		\draw[->,thick] (memory) -- (retrieve);
		\draw[->,thick] (retrieve) -- (motion);

	\end{tikzpicture}

	\vspace{0.5cm}

	\begin{alertblock}{Hypothesis}
		The droplet deposits a chemical trail that stores information about previously visited locations. By sensing gradients in this field, the droplet's future motion becomes influenced by its past motion.
	\end{alertblock}

\end{frame}

%------------------------------------------------
\begin{frame}{Research Question}

	\begin{center}

		\Large
		Does the observed self-avoidance indicate memory-driven dynamics?

	\end{center}

	\vspace{0.5cm}

	\begin{block}{Questions}
		\begin{enumerate}
			\item Does the droplet retain information about its past motion?
			\item Can this memory be quantified from trajectory data?
			\item Can a memory-driven model reproduce the observed behaviour?
		\end{enumerate}
	\end{block}

\end{frame}

%------------------------------------------------
\begin{frame}{How Do We Quantify Memory?}

	\centering

	\begin{tikzpicture}[
			node distance=2.2cm,
			every node/.style={
					draw,
					rounded corners,
					align=center,
					minimum width=3.2cm,
					minimum height=1cm
				},
			>=Stealth
		]

		\node (traj) {Trajectory Data};

		\node (stats) [right=of traj]
		{Statistical Analysis};

		\node (memory) [right=of stats,
			fill=red!25,
			draw=red!80!black,
			very thick]
		{Memory Quantification};

		% \node (memory) [right=of stats,
		% fill=red!15]
		% {Memory Quantification};

		\draw[->,thick] (traj) -- (stats);
		\draw[->,thick] (stats) -- (memory);

	\end{tikzpicture}

	\vspace{0.9cm}

	\begin{center}
		\large
		From observed trajectories, we extract statistical signatures of memory.
	\end{center}
	\begin{columns}
		\column{0.4\textwidth}

		\begin{block}{Transport Characterization}
			\begin{itemize}
				\item MSD
				      {\small \\How far does the droplet travel?}
				      % \item MSD \\- How far does the droplet travel?
				\item Local MSD Exponent
					      {\small \\Does transport change with time?}
			\end{itemize}
		\end{block}

		\column{0.4\textwidth}

		\begin{alertblock}{Memory Characterization}
			% \begin{block}{Memory Characterization}
			\begin{itemize}
				\item Velocity Autocorrelation
					      {\small \\How long is velocity remembered?}

				\item Orientation Correlation
					      {\small \\How long is direction remembered?}

			\end{itemize}
		\end{alertblock}

	\end{columns}


\end{frame}
